% ============================================================
\chapter{Mesure de Lebesgue sur $\R^n$}
\label{ch:lebesgue}
% ============================================================

Nous arrivons enfin \`a l'objet que toute cette machinerie pr\'epare~: la \emph{mesure de Lebesgue}. C'est elle qui donne un sens rigoureux \`a la notion de longueur, d'aire et de volume, l\`a o\`u la construction de Riemann \'echouait. Henri Lebesgue, dans sa th\`ese de 1902, a montr\'e que l'on peut mesurer des ensembles bien plus pathologiques que les simples r\'eunions d'intervalles --- et que l'int\'egrale qui en d\'ecoule poss\`ede des propri\'et\'es de passage \`a la limite incomparablement sup\'erieures \`a celle de Riemann. La construction que nous pr\'esentons ici applique le th\'eor\`eme de Carath\'eodory au contenu de volume des pav\'es, et produit la mesure la plus fondamentale de toute l'analyse.

\section{Construction}

La mesure de Lebesgue est l'unique mesure sur $\R^n$ qui \'etend la notion de volume
des pav\'es. Nous la construisons via le th\'eor\`eme de Carath\'eodory.

\begin{definition}[Pav\'e et volume]
Un \emph{pav\'e} (ou \emph{rectangle}) de $\R^n$ est un ensemble de la forme
$I = \prod_{k=1}^n [a_k, b_k)$ avec $a_k \leq b_k$. Son \emph{volume} est~:
\[
\operatorname{vol}(I) = \prod_{k=1}^n (b_k - a_k).
\]
\end{definition}

\begin{definition}[Alg\`ebre des r\'eunions de pav\'es]
Notons $\mathcal{A}_0$ la collection des r\'eunions finies disjointes de pav\'es
semi-ouverts de $\R^n$. C'est une alg\`ebre.
\end{definition}

\begin{proposition}
L'application $\mu_0 : \mathcal{A}_0 \to [0, +\infty]$ d\'efinie par
$\mu_0\!\left(\bigsqcup_{k=1}^m I_k\right) = \sum_{k=1}^m \operatorname{vol}(I_k)$
est une pr\'e-mesure $\sigma$-finie sur $\mathcal{A}_0$.
\end{proposition}

\begin{proof}
L'additivit\'e finie est imm\'ediate par d\'efinition. La $\sigma$-additivit\'e
d\'ecoule de la continuit\'e en $\varnothing$ (proposition~\ref{prop:critere_sigma_add}).
Soit $(A_m)$ une suite d\'ecroissante dans $\mathcal{A}_0$ avec $A_m \downarrow
\varnothing$. On montre que $\mu_0(A_m) \to 0$ par un argument de compacit\'e.

Supposons par l'absurde que $\mu_0(A_m) \geq \alpha > 0$ pour tout $m$. Chaque
$A_m$ est une r\'eunion finie de pav\'es. En r\'etr\'ecissant l\'eg\`erement chaque
pav\'e $[a_k, b_k)$ en $[a_k + \varepsilon, b_k - \varepsilon]$, on obtient un
compact $K_m \subset A_m$ avec $\mu_0(A_m) - \mu_0(K_m)$ petit. On peut s'assurer
que $\mu_0(K_m) \geq \alpha/2 > 0$, donc $K_m \neq \varnothing$. Comme
$K_1 \supset K_2 \supset \ldots$ sont des compacts non vides embo\^it\'es,
$\bigcap_m K_m \neq \varnothing$, ce qui contredit $\bigcap_m A_m = \varnothing$.

La $\sigma$-finitude vient de $\R^n = \bigcup_{k=1}^{\infty} [-k, k)^n$.
\end{proof}

\begin{theoreme}[Mesure de Lebesgue]
\label{thm:mesure_lebesgue}
Il existe une unique mesure $\lambda^n$ sur $(\R^n, \mathcal{B}(\R^n))$ telle que~:
\[
\lambda^n\!\left(\prod_{k=1}^n [a_k, b_k)\right) = \prod_{k=1}^n (b_k - a_k).
\]
Cette mesure est appel\'ee la \emph{mesure de Lebesgue} sur $\R^n$.
\end{theoreme}

\begin{proof}
Application directe du th\'eor\`eme d'extension de Carath\'eodory
(th\'eor\`eme~\ref{thm:extension_caratheodory}) \`a la pr\'e-mesure $\mu_0$. L'unicit\'e
vient de la $\sigma$-finitude.
\end{proof}

\begin{remarque}
On note souvent $\lambda = \lambda^1$ la mesure de Lebesgue sur $\R$.
La \emph{mesure de Lebesgue compl\`ete} est la compl\'etion de $\lambda^n$ par
rapport \`a la $\sigma$-alg\`ebre de Lebesgue $\mathcal{L}^n \supset
\mathcal{B}(\R^n)$.
\end{remarque}

\section{Propri\'et\'es fondamentales}

\begin{theoreme}[R\'egularit\'e]
\label{thm:regularite_lebesgue}
Pour tout $A \in \mathcal{B}(\R^n)$~:
\begin{enumerate}
    \item (R\'egularit\'e ext\'erieure) $\lambda^n(A) = \inf\{\lambda^n(U) :
    U \supset A,\, U \text{ ouvert}\}$.
    \item (R\'egularit\'e int\'erieure) $\lambda^n(A) = \sup\{\lambda^n(K) :
    K \subset A,\, K \text{ compact}\}$.
\end{enumerate}
\end{theoreme}

\begin{proof}
(1) Soit $\varepsilon > 0$. Par d\'efinition de la mesure ext\'erieure, il existe
$(I_k)$ suite de pav\'es avec $A \subset \bigcup_k I_k$ et
$\sum_k \operatorname{vol}(I_k) \leq \lambda^n(A) + \varepsilon$. En rempla\c{c}ant
chaque $[a_j, b_j)$ par $(a_j - \delta, b_j + \delta)$ avec $\delta$ assez petit,
on obtient un ouvert $U \supset A$ avec $\lambda^n(U) \leq \lambda^n(A) +
2\varepsilon$.

(2) Pour $\lambda^n(A) < \infty$, on utilise l'approximation
$A = \bigcup_m (A \cap [-m, m]^n)$ et la r\'egularit\'e ext\'erieure appliqu\'ee au
compl\'ementaire.
\end{proof}

\begin{theoreme}[Invariance par translation]
\label{thm:invariance_translation}
Pour tout $A \in \mathcal{B}(\R^n)$ et tout $x \in \R^n$~:
\[
\lambda^n(A + x) = \lambda^n(A), \quad \text{o\`u } A + x = \{a + x : a \in A\}.
\]
\end{theoreme}

\begin{proof}
La mesure $\nu(A) = \lambda^n(A + x)$ est une mesure sur $\mathcal{B}(\R^n)$
co\"incidant avec $\lambda^n$ sur les pav\'es (car le volume est invariant par
translation). Par unicit\'e, $\nu = \lambda^n$.
\end{proof}

\begin{theoreme}[Comportement par homoth\'etie]
Pour tout $A \in \mathcal{B}(\R^n)$ et tout $c \in \R \setminus \{0\}$~:
\[
\lambda^n(cA) = \abs{c}^n \lambda^n(A).
\]
Plus g\'en\'eralement, si $T : \R^n \to \R^n$ est une application lin\'eaire
inversible~:
\[
\lambda^n(T(A)) = \abs{\det T} \cdot \lambda^n(A).
\]
\end{theoreme}

\begin{proof}
Pour l'homoth\'etie, la mesure $\nu(A) = \lambda^n(cA)/\abs{c}^n$ co\"incide avec
$\lambda^n$ sur les pav\'es et est $\sigma$-finie, donc $\nu = \lambda^n$ par unicit\'e.

Pour le cas g\'en\'eral, on se ram\`ene \`a des transvections et des homoth\'eties
(d\'ecomposition de Gauss).
\end{proof}

\section{Ensembles de mesure nulle}

\begin{proposition}
Les ensembles suivants sont de mesure de Lebesgue nulle~:
\begin{enumerate}
    \item tout ensemble d\'enombrable\,;
    \item tout sous-espace affine de dimension $< n$ dans $\R^n$\,;
    \item l'ensemble de Cantor $\mathcal{C} \subset [0,1]$.
\end{enumerate}
\end{proposition}

\begin{proof}
(1) Un singleton $\{x\}$ est contenu dans $[x_1, x_1 + \varepsilon) \times \ldots
\times [x_n, x_n + \varepsilon)$ de volume $\varepsilon^n$, pour tout $\varepsilon > 0$.
Donc $\lambda^n(\{x\}) = 0$. Par $\sigma$-sous-additivit\'e, tout ensemble
d\'enombrable est de mesure nulle.

(2) Par invariance par les isom\'etries, il suffit de traiter le cas de l'hyperplan
$\R^{n-1} \times \{0\}$. On a $\R^{n-1} \times \{0\} \subset \R^{n-1} \times
[-\varepsilon, \varepsilon)$ pour tout $\varepsilon > 0$, et ce dernier ensemble
a mesure finie qui tend vers z\'ero.

(3) L'ensemble de Cantor est l'intersection $\mathcal{C} = \bigcap_n C_n$ o\`u
$C_n$ est une r\'eunion de $2^n$ intervalles de longueur $3^{-n}$. Donc
$\lambda(\mathcal{C}) \leq \lambda(C_n) = (2/3)^n \to 0$.
\end{proof}

\section{Ensembles non mesurables}

\begin{theoreme}[Vitali, 1905]
\label{thm:vitali}
Il existe un sous-ensemble de $[0, 1]$ qui n'est pas Lebesgue-mesurable.
\end{theoreme}

\begin{proof}
On d\'efinit la relation d'\'equivalence $x \sim y \iff x - y \in \Q$ sur $[0,1)$.
Par l'axiome du choix, soit $V \subset [0,1)$ un ensemble contenant exactement un
repr\'esentant de chaque classe d'\'equivalence.

Pour $r \in \Q \cap [0,1)$, d\'efinissons $V_r = \{x + r \pmod{1} : x \in V\}$
(translation modulo 1). Les $V_r$ sont deux \`a deux disjoints et
$[0,1) = \bigsqcup_{r \in \Q \cap [0,1)} V_r$.

Si $V$ \'etait mesurable, par invariance par translation (modulo 1), tous les $V_r$
auraient la m\^eme mesure $\alpha = \lambda(V)$. Par $\sigma$-additivit\'e~:
\[
1 = \lambda([0,1)) = \sum_{r \in \Q \cap [0,1)} \lambda(V_r) =
\sum_{r \in \Q \cap [0,1)} \alpha.
\]
Si $\alpha = 0$, la somme vaut $0 \neq 1$. Si $\alpha > 0$, la somme vaut
$+\infty \neq 1$. Contradiction.
\end{proof}

\begin{attention}[title=\textbf{R\^ole de l'axiome du choix}]
La construction de Vitali utilise l'axiome du choix. En effet, Solovay (1970) a
montr\'e qu'il est consistant avec ZF (sans l'axiome du choix) que tout
sous-ensemble de $\R$ soit Lebesgue-mesurable.
\end{attention}

\section{L'ensemble de Cantor}

\begin{definition}[Ensemble de Cantor]
L'ensemble de Cantor $\mathcal{C}$ est d\'efini par~:
\[
\mathcal{C} = \bigcap_{n=0}^{\infty} C_n, \quad \text{o\`u }
C_0 = [0,1], \quad C_{n+1} = \frac{C_n}{3} \cup \left(\frac{2}{3} + \frac{C_n}{3}\right).
\]
De mani\`ere \'equivalente, $\mathcal{C} = \{x \in [0,1] : x = \sum_{k=1}^{\infty}
a_k 3^{-k},\, a_k \in \{0, 2\}\}$.
\end{definition}

\begin{proposition}[Propri\'et\'es de l'ensemble de Cantor]
\begin{enumerate}
    \item $\mathcal{C}$ est compact, non d\'enombrable, parfait (ferm\'e sans point
    isol\'e), totalement discontinu.
    \item $\lambda(\mathcal{C}) = 0$.
    \item $\mathcal{C}$ est en bijection avec $[0,1]$ (via la fonction de Cantor).
    \item $\mathcal{C}$ contient des sous-ensembles non bor\'eliens (car
    $\abs{\mathcal{P}(\mathcal{C})} = 2^{\mathfrak{c}} > \mathfrak{c} =
    \abs{\mathcal{B}(\R)}$).
\end{enumerate}
\end{proposition}

\section{Caract\'erisation des ensembles Lebesgue-mesurables}

\begin{theoreme}
Un ensemble $A \subset \R^n$ est Lebesgue-mesurable si et seulement s'il existe un
bor\'elien $B$ et un ensemble de mesure nulle $N$ tels que $A = B \cup N$.
Autrement dit, $\mathcal{L}^n = \overline{\mathcal{B}(\R^n)}$ (la compl\'etion de
la tribu bor\'elienne).
\end{theoreme}

\begin{theoreme}[Caract\'erisation par les $G_\delta$ et $F_\sigma$]
Soit $A \subset \R^n$ Lebesgue-mesurable avec $\lambda^n(A) < \infty$. Alors~:
\begin{enumerate}
    \item il existe un $G_\delta$ (intersection d\'enombrable d'ouverts) $G \supset A$
    avec $\lambda^n(G \setminus A) = 0$\,;
    \item il existe un $F_\sigma$ (r\'eunion d\'enombrable de ferm\'es) $F \subset A$
    avec $\lambda^n(A \setminus F) = 0$.
\end{enumerate}
\end{theoreme}

\section{Exercices}

\begin{exercice}
Montrer que $\lambda(\Q) = 0$ directement par la d\'efinition de la mesure de
Lebesgue (en construisant un recouvrement explicite).
\end{exercice}

\begin{exercice}
Soit $A \subset \R$ mesurable avec $\lambda(A) > 0$. Montrer que l'ensemble des
diff\'erences $A - A = \{a - b : a, b \in A\}$ contient un intervalle centr\'e en
z\'ero. (\emph{Th\'eor\`eme de Steinhaus.})
\end{exercice}

\begin{exercice}[Ensemble de Cantor g\'en\'eralis\'e]
Pour $\alpha \in (0, 1)$, construire un ensemble ferm\'e $\mathcal{C}_\alpha
\subset [0, 1]$, hom\'eomorphe \`a l'ensemble de Cantor, avec
$\lambda(\mathcal{C}_\alpha) = 1 - \alpha$.
\end{exercice}

\begin{exercice}
Soit $f : \R \to \R$ lipschitzienne de constante $L$. Montrer que pour tout
$A \in \mathcal{B}(\R)$, $\lambda(f(A)) \leq L \cdot \lambda(A)$.
\end{exercice}

\begin{exercice}
Montrer que la mesure de Lebesgue $\lambda$ est la seule mesure (compl\`ete) sur
$\R$ qui est invariante par translation et telle que $\lambda([0,1]) = 1$.
\end{exercice}

\begin{exercice}
Montrer qu'il existe un ensemble Lebesgue-mesurable qui n'est pas bor\'elien.
\emph{Indication}~: utiliser la surjection $\mathcal{C} \to [0,1]$ et le fait
que $\mathcal{C}$ a mesure nulle.
\end{exercice}
